\documentclass[12pt]{article}
\usepackage[utf8]{inputenc}
\usepackage[T1]{fontenc}
\usepackage{hyperref}
\title{Enggano Language Resources for K-12 Curriculum Development}
\author{Compiled by Claude Sonnet 5.5 (claude-sonnet-5-5) for Andr\'as Kornai}
\date{}
\begin{document}
\maketitle
\section{Introduction}
Enggano (ISO 639-3 eno, Glottocode engg1245) is an endangered Austronesian language of Enggano island, Indonesia. An Oxford-led documentation project with partners including the University of Bengkulu is the main source of materials.
\section{Primary Educational Resources}
Search results state that the project aims to build educational materials for local schools, and a 2024 paper reports that Enggano schools currently do not use the local language. A PDF learner dictionary is mentioned on the Oxford project's pages.
\begin{itemize}
\item \url{https://www.humanities.ox.ac.uk/node/4320291}
\item \url{https://research.ox.ac.uk/map/enggano-in-the-austronesian-family-historical-and-typological-perspectives}
\end{itemize}
\section{Linguistic and Reference Materials}
\begin{itemize}
\item Enggano online dictionary (ACIS 2024 paper): \url{https://aisel.aisnet.org/acis2024/9}
\item Same paper, PDF: \url{https://enggano.ling-phil.ox.ac.uk/static/papers/20241016%20-%20ACIS%20Enggano%20Dictionary.pdf}
\item Enggano Digital Dictionary slides (Enggano, Indonesian, English, German): \url{https://enggano.ling-phil.ox.ac.uk/static/papers/20241206%20-%20Enggano%20Digital%20Dictionary.pdf}
\item Language Documentation \& Description journal article: \url{https://www.lddjournal.org/article/id/1009/}
\end{itemize}
\section{Religious and Community Resources}
None were found in this search.
\section{Digital and Audio Resources}
The project reports audio-visual recordings of conversations, stories and songs, and an online and print dictionary. Endangered Languages Project entries: \url{https://endangeredlanguages.com/node/38011/resources}, \url{https://endangeredlanguages.com/node/114524}
\section{Curriculum Development Context}
A 2026 article reports that official revitalization initiatives, including school curriculum integration, have largely fallen short: \url{https://jmic.online/issues/v15n1/16/}. I did not open these pages, so details should be checked.
\end{document}